\section{From attention to compute, memory, and latency}
\label{sec:fundamentals}
\subsection{What is stored in a KV cache?}
For a single attention head, let a token's representation be projected into a
query $q$, key $k$, and value $v$. Given a prefix of length $T$, causal attention
for its newest token has the form
\begin{equation}
 o_T=\operatorname{softmax}\!\left(\frac{q_T K_{1:T}^{\mathsf T}}{\sqrt{d_h}}\right)V_{1:T}.
 \label{eq:attention}
\end{equation}
The keys determine how strongly earlier positions are attended to; values
provide the content combined by those weights. This is the standard
scaled-dot-product attention construction~\cite{transformer}. The cache stores
previously computed keys and values at the required layers. It is neither the
model's learned weights nor a cache of finished textual answers.

For a causal model with an unchanged prefix, earlier hidden states do not depend
on future tokens. Their keys and values therefore need not be recomputed when
one token is appended. The new query must still attend to the appropriate old
state. Caching removes repeated prefix construction; it does not make reading
or attending over the prefix free. Queries normally need not be retained for
future autoregressive steps.

A cache entry has a computational identity. Matching text alone is insufficient
if model weights, adapters, tokenization, positional treatment, or multimodal
inputs differ. Reusing state under the wrong identity can silently change the
answer. Likewise, the KV for a passage following one prefix generally differs
from its KV following a different prefix. Exact prefix reuse and arbitrary
passage reuse are distinct mechanisms.

\subsection{Prefill and decode are different workloads}
During \emph{prefill}, the model processes known prompt tokens. Many token
projections and feed-forward operations can execute together as matrix-matrix
operations. During ordinary \emph{decode}, the next token is unknown until the
previous step completes. Each active sequence contributes roughly one new token
to an iteration, although speculative methods alter this unit of work.

Let $L$ be layers, $T$ prompt tokens, $d$ hidden width, and $H_qd_h=d$ for an
ordinary attention configuration. Ignoring architecture-dependent constants,
a dense model's prefill work has the form
\begin{equation}
 F_{\mathrm{prefill}}=\Theta(LTd^2)+\Theta(LT^2d).
 \label{eq:prefill}
\end{equation}
The first term collects projections and feed-forward computation, assuming its
intermediate width scales with $d$; the second is dense attention. A cached
single-token decode step at context length $t$ has work
\begin{equation}
 F_{\mathrm{step}}=\Theta(Ld^2)+\Theta(Ltd).
\end{equation}
For $U$ new output tokens after a $T$-token prompt, the attention contribution
sums as
\begin{equation}
 \sum_{u=0}^{U-1}\Theta(L(T+u)d)
 =\Theta\!\left(Ld\left(UT+\frac{U(U-1)}{2}\right)\right).
 \label{eq:decode}
\end{equation}
These expressions explain why a single cached decode step is linear in context
length while a long generation still accumulates a quadratic output-length
term. A deliberately naive implementation that reruns full dense attention over
the entire growing prefix at every step can accumulate cubic work in final
sequence length. That is not how a normal cached inference engine should run.
They are arithmetic counts, not wall-clock predictions: batching, memory
traffic, collective communication, and kernel utilization matter.

FlashAttention tiles exact attention to reduce traffic to high-bandwidth memory
and avoid materializing the full attention matrix there. It does not turn dense
attention into a linear-arithmetic algorithm~\cite{flash}. FlashAttention-3
uses hardware-oriented overlap and low-precision techniques~\cite{flash3}.
FlashInfer focuses on flexible, efficient inference kernels across cache
layouts and dynamic serving patterns~\cite{flashinfer}. These optimizations are
complementary to scheduling: a better admission policy cannot recover all the
performance lost by an unsuitable attention backend.

\subsection{The cache-size equation}
For a uniform full-attention model, an unsharded sequence's KV payload is
\begin{equation}
 M_{\mathrm{KV}}(T)=2LH_{kv}d_hbT,
 \qquad m_{\mathrm{token}}=2LH_{kv}d_hb,
 \label{eq:kvbytes}
\end{equation}
where $H_{kv}$ is the number of key/value heads, $b$ bytes per stored scalar,
and the leading two counts keys and values. This excludes allocator metadata,
quantization scales, temporary buffers, and replication. With $B$ independent
sequences of lengths $T_i$, payload is $m_{\mathrm{token}}\sum_iT_i$ before
sharing or architectural exceptions.

Multi-head attention (MHA) uses separate KV heads for query heads. Multi-query
attention (MQA) shares one KV head across queries~\cite{mqa}. Grouped-query
attention (GQA) sits between them, with several query heads sharing a KV
head~\cite{gqa}. Reducing $H_{kv}$ changes model architecture and its stored
representation; it is not merely a runtime eviction policy.

\paragraph{Worked example.}
Take an illustrative $L=32$, $H_{kv}=8$, $d_h=128$, and $b=2$. Then
\begin{align}
 m_{\mathrm{token}}&=2\cdot32\cdot8\cdot128\cdot2
 =131{,}072\ \mathrm{bytes}=128\ \mathrm{KiB},\\
 M_{\mathrm{KV}}(8192)&=1\ \mathrm{GiB}.
\end{align}
Sixty-four independent 8192-token contexts consume 64 GiB of payload before
weights and other allocations. A hypothetical 80 GiB device with 20 GiB of
weights and 8 GiB reserved for runtime overhead leaves 52 GiB, so only 52 such
contexts fit under this simplified arithmetic. This is a memory upper bound,
not a recommended concurrency or measured GPU capacity. Decode throughput and
latency may impose a lower limit.

\fig{kv-growth}{0.91}{Derived KV payload for an illustrative uniform model.
Changing the number of KV heads changes the slope; paging changes allocation
waste rather than the payload slope. No hardware was benchmarked.}{kvgrowth}

\subsection{Where the simple formula stops applying}
Multi-head latent attention, introduced in DeepSeek-V2, stores a compressed latent
representation rather than the ordinary full per-head KV tensors~\cite{mla}.
Its memory formula must include that representation and any positional state;
substituting an arbitrary small $H_{kv}$ into Equation~\ref{eq:kvbytes} is not a
general MLA estimator. Sliding-window layers retain a bounded attention window;
hybrid models may mix full attention, local attention, and recurrent state.
vLLM's hybrid cache manager explicitly handles different layer groups and their
cache-hit rules~\cite{hybrid}.

Tensor parallelism can shard cache heads, but heads or other state may be
replicated depending on the model and parallel degree. Pipeline parallelism
partitions layers. Context parallelism partitions sequence work/state and
introduces communication. Always distinguish total cluster bytes from bytes on
the busiest rank. Dividing a global KV estimate by GPU count is valid only if
the actual layout shards it evenly. A controller should consume measured
per-worker capacity and layout metadata rather than infer them from a model's
parameter count.

\subsection{The roofline view: what limits one iteration?}
A useful lower-bound model is
\begin{equation}
 t_{\mathrm{iter}}\gtrsim
 \max\left(\frac{F}{P_{\mathrm{eff}}},
           \frac{D_{\mathrm{HBM}}}{B_{\mathrm{HBM,eff}}},
           \frac{D_{\mathrm{link}}}{B_{\mathrm{link,eff}}}\right).
 \label{eq:roofline}
\end{equation}
$F$ is floating-point work; $D$ terms are transferred bytes; $P$ and $B$ are
achievable rates for this workload, not device marketing peaks. Real latency
also includes launches, synchronization, queueing, and portions that cannot
overlap. Adding all three ratios overestimates fully overlapped work; taking
only their maximum underestimates serial dependencies. The execution trace
determines which is appropriate.

Small-batch dense decode often pays heavily to read weights and KV for few new
tokens. Batching amortizes weight traffic across more tokens, but adds KV reads
and can increase per-iteration latency. Large prefill can exploit compute well,
yet long-context attention or distributed communication may become limiting.
``Prefill is compute-bound and decode is memory-bound'' is a useful first
hypothesis, not a property true for every shape, model, or device.

Arithmetic intensity is $I=F/D_{\mathrm{HBM}}$. In a simplified roofline,
achievable FLOPs/s is bounded by $\min(P_{\mathrm{eff}}, I B_{\mathrm{HBM,eff}})$.
Optimizing Python dispatch will have little effect while a long GPU kernel is
dominant and well fed. It can matter enormously when short iterations wait for
CPU scheduling or when a control service saturates a core. These are different
critical paths and require different measurements.

\subsection{Latency metrics and useful throughput}
Let $a_i$ be request arrival, $f_i$ first-token delivery, $z_i$ last-token
delivery, and $U_i$ the number of output tokens. Then
\begin{align}
 \mathrm{TTFT}_i &= f_i-a_i,\\
 \mathrm{TPOT}_i &= \frac{z_i-f_i}{U_i-1}\quad(U_i>1).
\end{align}
Individual inter-token latencies (ITLs) are adjacent delivery gaps. TPOT is an
average over those gaps and can hide a single long stall. TTFT includes queueing,
frontend work, cache retrieval, prefill, and delivery; it is not a direct prefill
kernel timer. Client-observed latency also includes transport and buffering.

For horizon $H$, define request goodput under a chosen joint SLO as
\begin{equation}
 G=\frac{1}{H}\sum_i\ind\{\mathrm{success}_i,
 \mathrm{TTFT}_i\le\tau_f,\ \mathrm{TPOT}_i\le\tau_o\}.
 \label{eq:goodput}
\end{equation}
Report the definition of success, errors, offered load, and rejected requests.
A system that rejects difficult work can appear to have excellent latency while
serving few users. Token throughput alone can favor longer outputs or a changed
workload mix. For agents, whole-job completion time and task success are needed
in addition to per-call metrics.
